% Part: sets-functions-relations
% Chapter: relations
% Section: special-properties

\documentclass[../../../include/open-logic-section]{subfiles}

\begin{document}

\olfileid{sfr}{rel}{prp}
\olsection{संबंधांचे विशेष गुणधर्म}

\begin{intro}
काही प्रकारचे संबंध इतके वारंवार आढळतात की त्यांना विशेष नावे दिली
आहेत. उदाहरणार्थ, $\le$ आणि $\subseteq$ हे आपापल्या प्रांतांतील घटकांना
सारख्याच प्रकारे जोडतात (उदा., $\le$ साठी $\Nat$ आणि $\subseteq$ साठी
$\Pow{A}$). त्यांच्यातील साम्य आणि भेद नेमके समजण्यासाठी संबंधांचे काही
विशेष गुणधर्मांनुसार वर्गीकरण करू. यांपैकी काही गुणधर्मांची संयोजने विशेष
महत्त्वाची ठरतात: क्रम आणि तुल्यता संबंध.
\end{intro}

\begin{defn}[परावर्तिता]
$R \subseteq A^2$ हा संबंध \emph{परावर्ती} तेव्हा आणि केवळ तेव्हाच असतो,
जेव्हा प्रत्येक $x \in A$ साठी $Rxx$ असते.
\end{defn}

\begin{defn}[संक्रमकता]
$R \subseteq A^2$ हा संबंध \emph{संक्रमक} तेव्हा आणि केवळ तेव्हाच असतो,
जेव्हा $Rxy$ आणि $Ryz$ असले, की $Rxz$ सुद्धा असते.
\end{defn}

\begin{defn}[सममिती]
$R \subseteq A^2$ हा संबंध \emph{सममित} तेव्हा आणि केवळ तेव्हाच असतो,
जेव्हा $Rxy$ असले, की $Ryx$ सुद्धा असते.
\end{defn}

\begin{defn}[प्रतिसममिती]
$R \subseteq A^2$ हा संबंध \emph{प्रति\-सम\-मित} तेव्हा आणि केवळ तेव्हाच
असतो, जेव्हा $Rxy$ आणि $Ryx$ दोन्ही असले, की $x=y$ असते (दुसऱ्या शब्दांत:
$x\neq y$ असेल, तर $\lnot Rxy$ किंवा $\lnot Ryx$ असते).
\end{defn}

\begin{explain}
सममित संबंधात $Rxy$ आणि $Ryx$ नेहमी एकत्र सत्य असतात किंवा दोन्ही असत्य
असतात. प्रतिसममित संबंधात $Rxy$ आणि $Ryx$ दोन्ही सत्य असण्यासाठी $x = y$
असणे आवश्यक असते. पण $x = y$ असताना $Rxy$ आणि $Ryx$ सत्य \emph{असलेच
पाहिजेत}, अशी अट यात नाही; फक्त त्यांना मनाई नाही. म्हणून प्रतिसममित
संबंध परावर्ती असू शकतो, पण प्रत्येक प्रतिसममित संबंध परावर्ती असतोच
असे नाही. तसेच, प्रतिसममित असणे आणि केवळ सममित नसणे या वेगळ्या अटी आहेत.
खरे तर, एकच संबंध एकाच वेळी सममित आणि प्रतिसममितही असू शकतो
(उदा., एकरूपता संबंध).
\end{explain}

\begin{defn}[संयोजितता]
$R \subseteq A^2$ या संबंधाला \emph{संयोजित} म्हणतात, जर सर्व $x,y\in A$
साठी $x \neq y$ असल्यास $Rxy$ किंवा $Ryx$ असते.
\end{defn}

\begin{prob}
पुढील गुणधर्म असणाऱ्या संबंधांची उदाहरणे द्या: (a) परावर्ती आणि सममित,
पण संक्रमक नाही; (b) परावर्ती आणि प्रतिसममित; (c) प्रतिसममित व संक्रमक,
पण परावर्ती नाही; आणि (d) परावर्ती, सममित व संक्रमक. संख्या किंवा संच
यांवरील संबंध वापरू नका.
\end{prob}

\begin{defn}[अपरावर्तिता]
$R \subseteq A^2$ या संबंधाला \emph{अपरावर्ती} म्हणतात, जर सर्व $x \in A$
साठी $Rxx$ असत्य असते.
\end{defn}

\begin{defn}[असममिती]
$R \subseteq A^2$ या संबंधाला \emph{असममित} म्हणतात, जर $x,y\in A$ अशा
कोणत्याही जोडीसाठी $Rxy$ आणि $Ryx$ दोन्ही सत्य नसतात.
\end{defn}

लक्षात घ्या: $A \neq \emptyset$ असेल, तर $A$ वरील कोणताही अपरावर्ती
संबंध परावर्ती नसतो, आणि $A$ वरील प्रत्येक असममित संबंध प्रतिसममितही
असतो. तरीही, काही $R \subseteq A^2$ संबंध परावर्तीही नसतात आणि
अपरावर्तीही नसतात; तसेच काही प्रतिसममित संबंध असममित नसतात.

\end{document}
